\section{Collapse phenomenology (untiered fleets)}
\label{sec:collapse}

\subsection{Anatomy of an absorbing collapse}
\label{sec:anatomy}

After the standard surge on an 8-replica untiered fleet at base rate
0.022 sessions/s, the system does not return to its pre-surge
operating point. It enters a cold state with hit rate $h$ 0.016--0.030
(pre-surge warm $h \sim 0.94$), TTFT p50 rising to 240\,s, queue growth
$\sim$1.1 requests/min to a waiting depth of 289, and throughput pinned
at 1.2--1.4 requests/s - the cold-branch service bound $\mu(0)$ of
Section~\ref{sec:occupancy}. No recovery occurs over the observed post-surge
horizons ($n=3$ runs; $\sim$73 post-surge minutes on the two 180-min
horizons, 195 on the 300-min horizon) \evid{R1}\evid{C01}. The state is
self-sustaining: the surge process no longer exists after
cancellation, and the base load that the same fleet served warm
before the surge maintains the collapse indefinitely on the measured
horizon.

The mechanism visible in the measured trajectories is the
drain-versus-catch-up race of Section~\ref{sec:occupancy}. The surge evicts the base
sessions' prefixes. After cancellation, the deferred base reservoir
drains back as catch-up flux; if that flux exceeds the cold-branch
service capacity before prefixes re-warm, misses regenerate the
eviction pressure and the fleet locks into the cold state \evid{R1}\evid{R2}.

\begin{figure}[t]
\centering
\includegraphics[width=\columnwidth]{fig_arcs_composite.png}
\caption{Measured trajectories, one per regime class, under the
identical eviction-scale surge (shaded, $t{=}62$--107\,min); all curves
are 5-min windows from per-replica scrapes, no model output. Warm
recovery at ($N{=}8$, 0.017): $h$ returns to $\sim$0.95, waiting $\sim$0. Absorbing
collapse at ($N{=}8$, 0.022) untiered: $h$ 0.013--0.030 after minute 200,
waiting climbing $\sim$1.4/min to $\sim$290 by minute 290. Delayed collapse at
($N{=}8$, 0.020) untiered: warm until $h$ falls from 0.89 at minute 240 to
0.05 by minute 280. Restore-bound congestion at ($N{=}8$, 0.022) with
the size-500 tier: $h$ 0.65--0.77 with tier-hit fraction 0.62--0.75
(dashed) and linear waiting growth $\sim$0.8/min. Complete recovery at
($N{=}8$, 0.022) with the size-375 tier: $h$ 0.97, waiting $\sim$0 at minute
300.}
\label{fig:arcs}
\end{figure}

\subsection{The collapse-probability band}
\label{sec:band}

The identical perturbation applied across base rates yields a
probability band, not a threshold line. Outcome counts over 17 runs
\evid{R2}\evid{C02}:

\begin{center}
\small
\begin{tabular}{@{}llp{0.52\columnwidth}@{}}
\toprule
$\ls$ (sessions/s) & runs & outcome \\
\midrule
0.012 & 2 & recovers 2/2 \\
0.017 & 4 & recovers 4/4 (150--180-min horizons; post-surge KV occupancy 0.40--0.78 across runs) \\
0.020 & 8 & three outcome classes: recovery $\times$3, delayed collapse $\times$2 (cold by $\sim$min 240), immediate collapse $\times$3 (onset min $\sim$130--150); collapse 5/8 by 300\,min \\
0.022 & 3 & collapses 3/3 \\
\bottomrule
\end{tabular}
\end{center}

Up to these $n$, the counts place the deterministic-collapse edge in
(0.020, 0.022]. These are outcome counts; no continuous
$p(\text{collapse} \mid \text{rate})$ curve is fitted to them. Among the 0.020 runs,
realized request rate does not order outcomes: one immediate
collapse occurred at realized 1.62 requests/s while one recovery
held at 1.78 \evid{R2}\evid{C04}. The band sits where Section~\ref{sec:breakpoints} places it:
above the cold-feasibility bound~\eqref{eq:coldbound} at 0.0098 sessions/s, well
below the warm-capacity limit, in the region where both branches are
locally stable and entry is fluctuation-decided.

\begin{figure}[t]
\centering
\includegraphics[width=\columnwidth]{fig_band_tally.png}
\caption{Outcome counts under the identical perturbation on the
8-replica untiered fleet at four base rates ($n=17$ runs): 0.012
recovers 2/2 and 0.017 recovers 4/4; 0.020 is bimodal over 8 runs
with recovery $\times$3, delayed collapse $\times$2, and immediate collapse $\times$3
(collapse 5/8 by 300\,min); 0.022 collapses 3/3, placing the
deterministic-collapse edge in (0.020, 0.022].}
\label{fig:bandtally}
\end{figure}

Two threshold axes are each indicated by one measured contrast pair
\evid{R2}\evid{C03}. On the duration axis, at 0.022 the 900\,s surge recovered
while the 2700\,s surge collapsed - consistent with the
reservoir-depth reading of Section~\ref{sec:breakpoints} (both surges fully evict the
pool; the longer one accumulates a six-fold deeper deferred
reservoir); a single 900\,s recovery is also consistent with chance
under any collapse probability materially below one. On the flux
axis, a rate-0.017 run served 4.9 requests/s of post-cancellation
catch-up traffic and stayed warm, while a rate-0.022 run re-entered
collapse at 4.8: the instantaneous catch-up rate does not
discriminate outcomes; the sustained arrival flux $\ls$ against
drain headroom does. No response curve is claimed on either axis.

\subsection{Horizon dependence, delayed collapse, and hysteresis}
\label{sec:horizon}

Surviving the perturbation does not imply long-term stability. At
($N{=}8$, 0.020), 2 of 8 runs that survived the surge and re-warmed
collapsed near minute 240 with no new perturbation: session
deepening gradually raises KV occupancy until it pins, and
per-replica pinning is near-synchronous (2 of 8 replicas to 8 of 8
within one 5-min window) \evid{R3}\evid{C05}. This delayed class is visible
only at the 300-min horizon; 60--180-min runs overstate stability at
this operating point, and the 0.017 row of the band table carries
the corresponding caveat (its 150--180-min horizons would not see a
delayed tail if one exists). In the theory's terms, the warm
equilibrium at this operating point sits close enough to the
separatrix that endogenous occupancy drift can cross it without any
exogenous trigger.

The cold state exhibits hysteresis. The same fleet that serves the
base load warm cannot re-warm from cold at that load, and recovery
requires dropping demand below the band's lower edge; in the
closed-loop experiments, each in-flight level has a unique
equilibrium and recovery occurs when the in-flight credit is lowered
\evid{R1}\evid{R8}\evid{C06}. The collapse is therefore a coexistence-band
phenomenon, not a capacity shortfall: the warm and cold branches
coexist at the same demand, as the theory's asymmetry argument
predicts (Section~\ref{sec:breakpoints}) - eviction runs on the pool-overwrite
timescale while re-warming runs on real time.

\subsection{The closed-loop contrast}
\label{sec:closedloop}

Closed-loop (completion-gated) arrivals self-stabilize instead of
collapsing. The measured boundary sits at achieved in-flight $\sim$73
warm versus $\sim$103+ degraded; degraded states are stationary rather
than absorbing; and the throughput function shows a flat segment
with cap-and-window co-movement (its falling segment is unreachable
statically and was not observed) \evid{R8}\evid{C07}. Absorbing collapse
requires demand that does not yield to backpressure: either
inelastic request-pinned generation or a deferred session reservoir
whose refill outpaces its drain \evid{R7}\evid{C07}. The open-loop
session-arrival protocol realizes the second condition - which is
also the condition agentic workloads realize in production. This
contrast locates the boundary between the queueing-stability
literature's assumptions (Section~\ref{sec:related}) and the regime where the
prefix-cache feedback dominates.

\subsection{The drain-depth separatrix}
\label{sec:separatrix}

The post-cancellation race has a measurable state variable: the
drain depth $D$, the fleet-mean KV-pool occupancy over minutes 115--120
of the protocol, the five-minute window opening eight minutes after
surge cancellation \evid{separatrix-findings.md}. $D$ is the empirical
coordinate of the separatrix of Section~\ref{sec:occupancy} (Figure~\ref{fig:fixedpoint}b): a shallow
drain (high $D$) means catch-up flux re-inflated the pool before
prefixes re-warmed, and the race is already lost when the window
closes.

At ($N{=}8$, 0.020), $D$ separates immediate collapse from all other
outcomes in all 8 runs (two of the five non-immediate runs later
collapse via the delayed mechanism): immediate-collapse minimum
0.735, other-outcome maximum 0.663 \evid{R15}\evid{C08}. The $D$ reading
precedes erosion onset in 20 of 21 fitted runs, and the same
occupancy read five minutes earlier does not separate (one survivor
read 0.674 mid-drain), so minute 115 is the earliest clean read
point.

Across the 21 fitted standard-surge runs (10 collapse events; one
16-replica run is excluded because its surge backlog was still
pinned at the read point and its horizon ends at 135\,min), a
Firth-penalized logistic fit of $p(\text{immediate} \mid D, \text{rate})$ gives $D$
signal beyond rate (penalized likelihood-ratio $p \sim 0.044$). The exact
stratified permutation test over the two mixed-outcome strata gives
one-sided $p = 3/112 = 0.027$. The signal is borne by the 8-replica
stratum: within ($N{=}8$, 0.020) the three immediate collapses hold
exactly the three highest $D$ values ($p = 1/56 = 0.018$), while the
$n=2$ 16-replica stratum at 0.037 is anti-aligned (the collapsing run
drained deeper than the surviving one) \evid{R15}\evid{C08}
\evid{out/separatrix\_draws.csv, separatrix\_fit.py}. The within-stratum
test treats the two dedicated-fleet runs and the six shard runs as
exchangeable; the two delayed collapses are exactly the two
dedicated-fleet runs, and the conservative shard-only form gives
$p = 1/20$. The fitted $p{=}0.5$ thresholds are $D^* = 0.79$, 0.69, 0.63 at
capacity-normalized rates 0.017, 0.020, 0.022; the 0.020 value sits
inside both the within-stratum separation interval (0.663, 0.735)
and the pooled 8-replica interval across all three rates. At $n=21$
the threshold's direction and location are established, its
steepness is not (the $D$ coefficient's Wald 95\% interval spans
roughly (1, 39)). $D$ also orders the collapse-onset spectrum: among
the 12 fitted collapsing runs, pooled across fleet sizes, rates, and
collapse classes (all of which co-vary with $D$ and onset), the rank
correlation between $D$ and onset time is $-0.99$ - shallower drains
collapse earlier. The covariate was identified mid-campaign, so
these $p$-values are not from a pre-registered test; the three runs
measured after identification landed on the hypothesized sides of
the threshold.

Two scope limits bound the claim. First, the threshold moves with
rate and fleet size: at $N{=}16$, the 0.034 point survived $D = 0.74$ and
a 0.037 run collapsed from 0.65, each a single-run point observation
\evid{R15}\evid{C09} - as expected if $D$ is one coordinate of a separatrix that
also depends on rate and span, not a universal constant. Second, $D$
does not predict the delayed class: those runs win the drain race
from deep drains ($D$ 0.616 and 0.663) and collapse $\sim$2\,h later by the
distinct occupancy-drift mechanism of Section~\ref{sec:horizon} \evid{R15}\evid{C09}. $D$ is a
mediating observable of the drain race, not a control law. In the
pooled fit it is confounded with the $N{=}16$ indicator at $n=21$, because
every 16-replica run at the 0.020-equivalent point both drained
shallow and collapsed; the mediation interpretation rests on the
within-stratum permutation test and the model comparison of
Section~\ref{sec:models} \evid{C08}.

\begin{figure}[t]
\centering
\includegraphics[width=\columnwidth]{separatrix_fit.png}
\caption{Drain-depth separatrix. $D$ (fleet-mean KV occupancy,
minutes 115--120) versus capacity-normalized rate for the 21 fitted
runs, marked by outcome class; the line is the Firth-penalized
$p(\text{immediate}) = 0.5$ contour ($D^* = 0.79/0.69/0.63$ at
0.017/0.020/0.022). The two delayed collapses sit at the deep end
and are not predicted by $D$; the 16-replica points illustrate the
threshold's span dependence.}
\label{fig:separatrix}
\end{figure}
